\begin{textsample}{POS Dim 1 – rr – Score 40.00 – rr/rr\_ef\_67.txt}  \label{ex:f1_pos_262}
5.6 AVALIAÇÃO DA APRENDIZAGEM A avaliação do processo ensino e aprendizagem tem sido \textbf{uma} preocupação constante dos professores, pois faz parte do trabalho docente verificar e julgar o rendimento dos alunos, avaliando os resultados do ensino, cabendo também ao professor reconhecer as diferenças na capacidade de aprender dos alunos, para ajudá -los a superar suas dificuldades e avançar na aprendizagem.
Os dados sobre avaliação da aprendizagem, em Roraima, de acordo com o último IDEB ( 2017 ), são, no mínimo, preocupantes.
De acordo com a metodologia utilizada, a Escala de Proficiência em Língua Portuguesa e Matemática no 9º ano do Ensino Fundamental e no 3º ano do Ensino Médio são os seguintes : Pelos dados acima é possível perceber \textbf{que} o componente de Língua Portuguesa não se apresenta como instrumento, na medida em \textbf{que} o aluno não \textbf{consegue} utilizá -la como recurso ou ferramenta para expressar -se no mundo.
Parece contraditório \textbf{que} no \textbf{século} XXI, caracterizado como a era da comunicação e a \textbf{democratização} do acesso à informação, os jovens com total domínio dos meios eletrônicos como WhatsApp, Twitter, SMS, Facebook etc., não \textbf{consigam} articular saberes necessários para a produção textual, usando a norma mais elaborada da língua escrita nem extrair ideias e interpretar \textbf{um} texto, mesmo aqueles de estrutura mais \textbf{simples}.
No componente Matemática, os dados são ainda mais alarmantes, conforme se \textbf{vê} na tabela seguinte.
Observa -se \textbf{que} os níveis de aprendizagem dos alunos dos anos finais das duas principais etapas da Educação Básica estão aquém do desejado, o \textbf{que} deve ser \textbf{visto} como desafios para os professores atuantes nesse nível de ensino, pois \textbf{um} dos papéis da avaliação é fazer com \textbf{que} o professor repense sua prática pedagógica a partir dos resultados obtidos.
Embora esses resultados possam expressar a interferência de múltiplos fatores, alguns precisam ser assumidos pelo docente para redimensionar suas atividades.
Nesse sentido, as 10 ( dez ) Competências Gerais da BNCC trazem alguns princípios \textbf{que} devem ser levados em conta no processo avaliativo. \textbf{Um} deles é a perspectiva de aprendizagem contínua, ou seja, apesar do currículo está organizado em ano e série, as aprendizagens não “ respeitam ” esse limite.
Elas são construídas contínua e cumulativamente, e aquilo \textbf{que} não foi alcançado em \textbf{um} momento poder perfeitamente ser desenvolvido em outro. \textbf{Historicamente}, a educação brasileira é marcada pelo processo autoritário de dominação, submissão e reprodução das desigualdades e, em consequência, a avaliação educacional segue essa mesma lógica.
A ideologia do sistema capitalista liberal tem suas origens nos princípios da competitividade, da eficiência, da reprodução e \textbf{exclusão} e, por \textbf{conseguinte}, se manifesta na prática educativa e avaliativa sob vários aspectos : \textbf{um} deles é quando se utiliza da avaliação para criar \textbf{um} ambiente escolar de competitividade por \textbf{um} espaço na sociedade.
Se de \textbf{um} \textbf{lado}, as \textbf{pedagogias} \textbf{que} fazem parte do pensamento liberal, como a tradicional, a progressivista e a tecnicista dão ênfase ao modelo de avaliação \textbf{que} prima pela repetição e reprodução, têm -se também as \textbf{pedagogias} não-diretivas \textbf{que} se centram no aluno e podem negar o saber elaborado por ser fruto do pensamento da classe dominante.
Cada \textbf{uma} das tendências avaliativas serve a \textbf{um} modelo de educação, manifestando -se na prática do professor de diferentes formas : o professor pode ser o detentor do saber e cabe ao aluno apenas ser o “ receptáculo ” de informações ; nesse modelo Foucault ( 2003 ) \textbf{argumenta} \textbf{que} a escola se parece com prisões, onde as regras se sobrepõem à didática e à relação de colaboração entre professor e aluno.
Ou, como nas \textbf{pedagogias} não-diretivas, a avaliação pode \textbf{perder} seu significado de retroalimentar o processo, porque se apresenta como excludente.
Assim, a avaliação no âmbito educacional deve ser entendida como ação \textbf{que} atenda de forma qualitativa o acompanhamento dos processos de ensino e aprendizagem e, como tal, requer \textbf{uma} prática pedagógica \textbf{que} proporcione informações relevantes sobre esses processos.
De acordo com vários \textbf{autores}, entre eles, Libâneo \textbf{et} al. ( 2008 ), as práticas avaliativas são meios para intervir nos problemas detectados, não servem como prêmio ou castigo para as instituições escolares, nem para os alunos.
A prática avaliativa de caráter formativo é capaz de proporcionar informações qualitativas sobre o desenvolvimento dos processos de ensino e de aprendizagem, e informar em \textbf{que} medida o aluno está se desenvolvendo, quais as competências foram formadas, etc.
Neste sentido, deve centrar -se em questões qualitativas \textbf{que} envolvam ações com vistas a \textbf{uma} tomada de decisão durante o processo.
Considerando a \textbf{abordagem} crítica da educação, a avaliação como processo formativo deve ter como objetivo a aprendizagem do aluno e para isto requer \textbf{uma} prática pedagógica voltada para a superação da avaliação como instrumento de dominação e \textbf{exclusão}.
A Lei de Diretrizes e Base da Educação Nacional 9.394/96, os Parâmetros Curriculares Nacionais, as Diretrizes Curriculares Nacionais para a Formação de Professores da Educação Básica, em Nível Superior ( Resolução CNE/CP 01/2002 ), defendem a avaliação como \textbf{uma} prática cidadã, alicerçada na construção do conhecimento, com qualidade formal e política \textbf{que} rompa com a concepção de caráter tradicional e excludente.
Os estudos \textbf{atuais} se debruçam sobre o tema avaliação concebendo -a como \textbf{um} conjunto de conhecimentos e habilidades essenciais e \textbf{imprescindíveis} à formação do professor e sua profissionalização docente e tem como foco a função reflexiva, crítica e investigativa sobre os processos de ensinar e aprender.
Nesta perspectiva, \textbf{salienta} -se a importância da prática pedagógica e avaliativa alicerçada na perspectiva de reflexão-ação-reflexão sobre os processos de ensino e de aprendizagem.
Não basta avaliar é \textbf{preciso} tomar decisões qualitativas sobre os resultados da avaliação institucional e pedagógica.
Assim, é necessária a compreensão de \textbf{que} ensinar não é transferir conhecimentos e, sim, constitui -se num vínculo \textbf{que} envolve professor e alunos num ato de acolhimento, doação, reciprocidade, respeito e ajuda, \textbf{que} nega a avaliação como instrumento de poder e \textbf{exclusão}.
A avaliação como \textbf{inerente} à atividade humana se faz presente no âmbito educacional como parte integrante do processo de ensino e aprendizagem, o \textbf{que} lhe outorga a \textbf{categoria} de componente essencial desse processo.
Nesta \textbf{direção}, defende -se a avaliação como prática social, \textbf{que} \textbf{remete} a \textbf{uma} reflexão dos processos de ensinar e aprender como \textbf{constituintes} de \textbf{uma} prática formativa por \textbf{excelência}.
As práticas avaliativas realizadas no interior das salas de aula não são \textbf{neutras}.
Elas corroboram com a concepção de educação adotada em cada ambiente escolar.
Na medida em \textbf{que} se utilizam de \textbf{provas} e exames para testar a aprendizagem como pensamento da classe dominante.
Segundo Vasconcellos ( 2007 ), romper com a prática da avaliação classificatória não é fácil, pois seu viés ideológico como instrumento de controle, de inculcação ideológica e de discriminação social está impregnado no ideário pedagógico da maioria dos professores.
Segundo Libâneo \textbf{et} al. ( 2008 ), as práticas avaliativas concebidas como processo são meios para intervir nos problemas de aprendizagem e servir de parâmetro para a melhoria do ensino.
Dessa forma, sustentam -se em parâmetros qualitativos referentes ao ensino ministrado e à aprendizagem adquirida.
Avaliar na concepção de \textbf{um} projeto educativo de qualidade é diferente de medir, testar, classificar e excluir.
Portanto, em contraposição a estes termos, considera -se a avaliação como meio de acompanhamento do processo de ensino e da aprendizagem, para \textbf{uma} tomada de decisão.
Quando se \textbf{argumenta} \textbf{que} avaliar \textbf{implica} \textbf{uma} ação \textbf{que} envolve julgamento com vistas a \textbf{uma} tomada de decisão, essa ação deve estar balizada por princípios, sobretudo éticos, e o julgamento ético \textbf{implica} em \textbf{um} movimento de “ dentro para fora ”, ou seja, não existe a ética dos outros, mas a minha ética, daquilo \textbf{que} acredito, daquilo \textbf{que} faz parte da moral individual \textbf{que} se manifestam no momento de julgar, \textbf{uma} vez \textbf{que} as decisões decorrentes da avaliação da aprendizagem \textbf{implicam} \textbf{encaminhamentos} na vida escolar dos estudantes, tais como reorientações de percurso ao longo do ano, as chamadas recuperações paralelas, exames finais e até mesmo a decisão acerca da reprovação escolar e suas implicações.
Além dessas existem outras consequências : quem continua e quem está fora da escola, quem tem condições de ingressar no mercado de trabalho e quem pode ou não fazer ensino superior. \textbf{Contudo}, a avaliação precisa ser \textbf{vista} como parte do processo educacional, mas deve ser repensada como mecanismo \textbf{que} inclui, \textbf{que} detecta falhas e \textbf{que} enaltece avanços.
Razão pela qual ela deve ser pensada junto a outros elementos formadores do currículo, como a metodologia utilizada, os recursos e a didática do professor.
Assim, a visão de \textbf{sujeito} \textbf{que} se quer formar, papel social da escola e o projeto de sociedade \textbf{que} se imagina a partir das políticas educacionais, está \textbf{implicada} e refletem o modelo ou concepção de avaliação \textbf{que} se tem.
Currículo e avaliação têm relações estreitas não só do ponto de vista \textbf{teórico}, mas também da prática.
Pois, diferentes \textbf{abordagens}, embora possam identificar -se mais ou menos com determinadas posturas \textbf{teóricas}, trazem no seu \textbf{bojo} questões históricas e sociais de determinada época.
Há \textbf{que} se considerar ainda a complexidade das ações educativas e as interferências do \textbf{sujeito} político, professor, no convívio diário da sala de aula na qual se materializa o currículo, \textbf{que}, como sabemos, não é \textbf{neutro} ou destituído de intencionalidade e sim \textbf{um} instrumento político e epistemológico.
Por essa razão, é possível dialogarem diferentes \textbf{abordagens} sobre avaliação num mesmo currículo, sem, \textbf{contudo}, conflitarem -se.
No plano prático é possível avaliar a aprendizagem sem avaliar o currículo, mas não é possível avaliar o currículo sem avaliar a aprendizagem, de maneira \textbf{que} quando se delineia \textbf{um} currículo é \textbf{preciso} ter definido o modelo e concepção de aprendizagem e suas formas de avaliar.
A partir do \textbf{que} a BNCC trouxe em relação às competências e habilidades, \textbf{que} se quer formar, é possível se delinear possíveis perspectivas de avaliação.
De acordo com esse documento, a sociedade \textbf{contemporânea} impõe \textbf{um} olhar inovador e inclusivo a questões \textbf{centrais} do processo educativo : o \textbf{que} aprender, para \textbf{que} aprender, como ensinar, como promover redes de aprendizagem colaborativa e como avaliar o aprendizado.
A perspectiva crítica nos parece a \textbf{que} melhor responde a esses interrogantes : os porquês, os para quê e como, cuja \textbf{centralidade} é o aluno.
Nessa perspectiva, as aprendizagens estão no centro do processo, em torno da qual gravitam outros elementos, dentre eles a avaliação, mas não aquela como sinônimo de verificações ou medições.
No cotidiano da sala de aula, podemos traduzir essa perspectiva como atividades mais participativas, ações relacionadas à construção da autonomia, aos processos de acompanhamento dos estudantes em suas múltiplas possibilidades, ao respeito aos diferentes ritmos e tempos de aprendizagem dos \textbf{sujeitos}.
No \textbf{que} tange às competências e às centenas de habilidades previstas para o componente de Língua Portuguesa, é possível vislumbrar \textbf{um} \textbf{sujeito} crítico, criativo, ético e competente, com apropriação do conhecimento e com capacidade para compreendê -lo, empregá -lo e recriá -lo.
Assim, a avaliação da aprendizagem desse aluno não se coaduna com questões positivistas, com múltipla escolha ou com reprodução do conhecimento.
Mas, com \textbf{uma} avaliação \textbf{que} ultrapasse a avaliação quantitativa e classificatória, sem necessariamente, dispensá -la.
A partir do \textbf{que} a BNCC propõe, destaca -se a importância da avaliação com caráter formativo, considerando sua importância em relação às competências e habilidades \textbf{que} se desenvolvem no decorrer do processo.
Assim, deve -se atender a formação de \textbf{um} \textbf{sujeito} crítico capaz de dialogar com os conhecimentos adquiridos para a superação de desafios impostos na sociedade \textbf{contemporânea}.
Portanto, a avaliação \textbf{que} dê conta de formar \textbf{um} aluno com as competências e habilidades previstas na BNCC, \textbf{certamente} ultrapassará as questões quantitativas e se apresentará de maneira a atender as necessidades dos alunos, nas quais possam emitir opiniões, refletir e apresentar proposições para soluções de conflitos presentes na sua realidade.

% matched lemmas: abordagem, argumentar, atual, autor, bojo, categoria, central, centralidade, certamente, conseguinte, conseguir, constituinte, contemporâneo, contudo, democratização, direção, encaminhamento, et, excelência, exclusão, historicamente, implicar, imprescindível, inerente, lado, neutro, pedagogia, perder, preciso, prova, que, remeter, salientar, simples, sujeito, século, teórico, um, uma, ver
\end{textsample}
